\section*{Supplemental Tables}

% Stats for the different stages of the assembly
\begin{table}[ht]
    \centering
    \suptablecaption{\textbf{Assembly statistics for the \balgla{} genome
        at different stages of the assembly process.}
        Default contigs refers to the sequences assembled by \textsc{hifiasm} with
            default parameters.
        Optimized contigs denotes the sequences generated using \textsc{hifiasm}
            followed by removal of contaminant sequences and purging of haplotigs.
        Hi-C scaffolds refer to the scaffolds generated by \textsc{YaHS}.
        Curated assembly refers to the sequences generated after manual curation of 
            the Hi-C contact map and polishing with \textsc{Inspector}.
        \textsc{BUSCO} results refer to the \texttt{arthropoda\_odb10} dataset.}
    \shadedrows
    \begin{tabular}{lrrrr}
        \toprule
        Assembly Stage &
            \shortstack{Default \\ contigs} &
            \shortstack{Optimized \\ contigs} &
            \shortstack{Hi-C \\ scaffolds} &
            \shortstack{Curated \\ assembly} \\
        \midrule \startshading
        Sequence length (Mbp)          & 1.64        & 1.05        & 1.05       & 1.04 \\
        Number of fragments            & \num{4212}  & \num{1342}  & 650        & 592 \\
        Contig N50 (Mbp)               & 0.734       & 1.35        & 1.18       & 1.18 \\
        Scaffold N50 (Mbp)             & \texttt{NA} & \texttt{NA} & 50.96      & 51.41 \\
        \textsc{BUSCO} complete (\%).  & 95.66       & 93.88       & 93.98      & 94.08 \\
        \textsc{BUSCO} single (\%)     & 45.31       & 89.04       & 89.24      & 89.44 \\
        \textsc{BUSCO} duplicates (\%) & 50.35       & 4.84        & 4.74       & 4.64 \\
        \textsc{BUSCO} fragmented (\%) & 0.79        & 1.18        & 0.89       & 0.89 \\
        \textsc{BUSCO} missing (\%)    & 3.55        & 4.94        & 5.13       & 5.03 \\
        \textsc{BUSCO} size            & \num{1013}  & \num{1013}  & \num{1013} & \num{1013} \\
        \bottomrule
    \end{tabular}
    \label{tab:sup_asm_stats}
\end{table}

% Gene-completeness in the annotation
\begin{table}[ht]
    \centering
    \suptablecaption{\textbf{Completeness metrics for the \balgla{} protein-coding
        gene annotation.}
        \textsc{BUSCO} completeness was determined by \textsc{compleasm}
        \texttt{protein}.}
    \shadedrows
    \begin{tabular}{lrr}
        \toprule
        Metric &
            \texttt{arthropoda\_odb10} &
            \texttt{crustacea\_odb12} \\
        \midrule \startshading
        \textsc{BUSCO} complete (\%)   & 93.58      & 84.77 \\
        \textsc{BUSCO} single (\%)     & 79.76      & 72.66 \\
        \textsc{BUSCO} duplicates (\%) & 13.82      & 12.11 \\
        \textsc{BUSCO} fragmented (\%) & 1.78       & 4.10 \\
        \textsc{BUSCO} missing (\%)    & 4.64       & 11.13 \\
        \textsc{BUSCO} size            & \num{1013} & \num{1536} \\
        \bottomrule
    \end{tabular}
    \label{tab:sup_annot_stats}
\end{table}

% Repeat annotation
\begin{table}[htbp]
\centering
\suptablecaption{\textbf{High-level repeat annotation of the \balgla{}
        assembly.}
            Repeats were first identified with the \textsc{EarlGrey} software.
            The classification is the merge of \textsc{EarlGrey}'s homology-based
            assignment, followed by the re-classification of unknown elements by
            \textsc{TEsorter} and \textsc{DeepTE}.
            Coverage is the total number of masked sites per class; ``\% Genome''
            is the coverage divided by assembly size (\num{1043519078}\,bp).}
    \label{tab:sup_repeat_summary}
    \shadedrows
    \begin{tabular}{lrrrr}
        \toprule
            Classification &
                Coverage (bp) &
                Count &
                \% Genome &
                \shortstack{Distinct \\ elements} \\
        \midrule \startshading
            DNA                                        & \num{379178613} & \num{895512}  & 36.34 & 35 \\
            LINE                                       & \num{107103442} & \num{223314}  & 10.26 & 20 \\
            LTR                                        & \num{50161694}  & \num{120887}  &  4.81 &  8 \\
            \shortstack[l]{Other (Simple Repeat, \\
                ~~Microsatellite, RNA)}                & \num{37359678}  & \num{74786}   &  3.58 &  3 \\
            Penelope                                   & \num{25833119}  & \num{65921}   &  2.48 &  4 \\
            Rolling Circle                             & \num{22238826}  & \num{42257}   &  2.13 &  1 \\
            Retroposon                                 & \num{2095328}   & \num{5834}    &  0.20 &  2 \\
            SINE                                       & \num{1253619}   & \num{4437}    &  0.12 &  3 \\
            DIRS                                       & \num{1184091}   & \num{1285}    &  0.11 &  2 \\
            Unclassified                               & \num{61450725}  & \num{148647}  &  5.89 &  1 \\
            \midrule
            Total                                      & \num{687859135} & \num{1582880} & 65.92 & 79 \\
    \bottomrule
    \end{tabular}
\end{table}

% Stats about polymorphism in the reference individual and the popgen samples.
\begin{table}[ht]
    \centering
    \suptablecaption{\textbf{Genetic variation in central Oregon \balgla{}.}
        Percentage of variant sites is calculated for single nucleotide
        variants ($100 \times N_\text{SNP} / N_\text{accesible}$).
        Density of SNPs per accessible kbp is calculated only for windows with more
        than 100 accessible sites ($1000 \times N_\text{window SNPs} / N_\text{window accessible}$).
        Density of SNP in coding sequences was calculated only for exons with more than
        100 accessible sites and adjusting for exon length
        ($1000 \times N_\text{exon SNPs} / N_\text{exon accessible}$).}
    \shadedrows
    \begin{tabular}{lrr}
        \toprule
        Metric &
            Reference (N=1) &
            Cape Perpetua (N=6) \\
        \midrule \startshading
        Reference sequence length (bp)        & \num{904940360} & \num{904940360} \\
        Accessible sites post-filtering (bp)  & \num{320031567} & \num{231589201} \\
        Total variants                        & \num{9587325}   & \num{21395217} \\
        Indels                                & \num{730793}    & \num{899926} \\
        SNPs                                  & \num{8856532}   & \num{20495291} \\
        Percent variant sites (\%)            & 2.77            & 8.85 \\
        Multiallelic SNPs                     & \texttt{NA}     & \num{910949} \\
        Percent multiallelic SNPs (\%)        & \texttt{NA}     & 4.6 \\
        \midrule
            \rowcolor{white}SNP density \\
        \midrule
        % This is 1000 x N_snps / N_accessible in 1000 bp windows
        Mean SNPs per accessible kbp          & 26.2           & 83.3 \\
        Median SNPs per accessible kbp        & 19.0           & 80.7 \\
        StDev SNPs per accessible kbp         & 27.2           & 53.9 \\
        \midrule
            \rowcolor{white}Coding variants (all) \\
        \midrule
        % This is 1000 x N_exon_snps x N_exon_accessible across each exon interval
        SNPs in CDS                           & \num{396681}    & \num{1061714} \\
        Percent SNPs in CDS (\%)              & 4.48            & 5.18 \\
        Mean SNPs per accessible CDS kbp      & 15.8            & 51.0 \\
        Median SNPs per accessible CDS kbp    & 12.0            & 46.5 \\
        StDev SNPs per accessible CDS kbp     & 17.0            & 33.8 \\
        \midrule
            \rowcolor{white}Coding variants (non-synonymous) \\
        \midrule
        % Just counting non-synomymous SNPs
        Non-syn. SNPs in CDS                  & \num{144620}    & \num{337040} \\
        Percent of coding SNPs (\%)           & 36.5            & 31.7 \\
        % Just counting per total gene, regardless of size
        Mean non-syn. SNPs per gene           & 8.2             & 20.0 \\
        Median non-syn. SNPs per gene         & 3.0             & 12.0 \\
        StDev non-syn. SNPs per gene          & 15.4            & 27.4 \\
        % Now adjusting for size. Total Non-syn SNP density
        Mean non-syn. SNP per accessible CDS kbp    & 7.9  & 21.6 \\
        Median non-syn. SNP per accessible CDS kbp  & 3.2  & 13.3 \\
        StDev non-syn. SNP per accessible CDS kbp   & 13.2 & 24.3 \\
        \bottomrule
    \end{tabular}
\label{tab:sup_popgen_stats}
\end{table}

% Stats about the phasCons conserved elements and enrichment of genomic elements.
\begin{sidewaystable}[ht]
    \suptablecaption{\textbf{Enrichment of annotated features among \textsc{phastCons}
        highly conserved elements.}
        The total feature proportion was determined according to the \num{841166098}\,bp
        corresponding to the 16 chromosome-level scaffolds.
        Proportion of highly-conserved elements (HCE) was calculated based on the
        \num{50849898}\,bp-span of all HCEs in the genome.
        The number of chromosomes denotes the number of chromosome-level scaffolds
        containing annotations of the given feature.}
    \centering
    \label{tab:sup_phastcons_enrichment}
    \shadedrows
    \begin{tabular}{lrrrrrrrr}
    \toprule
        Feature &
            \shortstack{Total feature \\ length (bp)} &
            \shortstack{Total feature \\ proportion} &
            \shortstack{HCE feature \\ length (bp)} &
            \shortstack{HCE feature \\ proportion} &
            \shortstack{Mean $\log_2$ \\ enrichment} &
            \shortstack{Median $\log_2$ \\ enrichment} &
            \shortstack{St Dev $\log_2$ \\ enrichment} &
            \shortstack{Number of \\ chromosomes} \\
    \midrule \startshading
        3' UTR     & \num{2358104}   & \num{2.80e-3} & \num{152598}   & \num{3.00e-3} & $-0.0861$    & 0.0332       & 0.508        & 16 \\
        5' UTR     & \num{530151}    & \num{6.30e-4} & \num{27679}    & \num{5.44e-4} & $-0.292$     & $-0.332$     & 0.453        & 16 \\
        CDS        & \num{16273411}  & 0.0193        & \num{2225138}  & 0.0438        & 0.752        & 0.984        & 0.926        & 16 \\
        TE         & \num{370573528} & 0.441         & \num{17505920} & 0.344         & $-0.359$     & $-0.305$     & 0.180        & 16 \\
        intergenic & \num{196612776} & 0.234         & \num{10688761} & 0.210         & $-0.167$     & $-0.178$     & 0.167        & 16 \\
        intron     & \num{111562654} & 0.133         & \num{8575960}  & 0.169         & 0.419        & 0.297        & 0.432        & 16 \\
        lncRNA     & \num{2211027}   & \num{2.63e-3} & \num{130744}   & \num{2.57e-3} & $-0.174$     & $-0.133$     & 0.567        & 16 \\
        multiple   & \num{141022121} & 0.168         & \num{11537936} & 0.227         & 0.418        & 0.438        & 0.167        & 16 \\
        rRNA       & \num{2106}      & \num{2.50e-6} & 0.000          & 0.000         & \texttt{NaN} & \texttt{NaN} & \texttt{NaN} & 6 \\
        tRNA       & \num{20220}     & \num{2.40e-5} & \num{5162}     & \num{1.02e-4} & 2.143        & 2.041        & 0.959        & 16 \\
   \bottomrule
\end{tabular}
\end{sidewaystable}

% Breakdown of per-species codon usage bias
\begin{sidewaystable}[ht]
    \centering
    \small
    \suptablecaption{\textbf{Codon usage biases across metazoan genomes.}
        \comment{update (remove?) this}
        Table shows the codon usage statistics (effective number of codons; ENC), GC
        proportion, and GC proportion at third codon positions (GC3s) for length-filtered
        (CDS $\geq$ 300\,bp, within 2 standard deviations of the per-species mean length)
        genes in each genome.
        The first four rows (\texttt{BalGla}, \texttt{BalCre}, \texttt{AmpAmp}, and
        \texttt{PolPol}) denote the focal barnacle genomes.}
    \label{tab:sup_enc_summary}
    \shadedrows
    \begin{tabular}{llrrrrrrrrrr}
    \toprule
    & & & \multicolumn{3}{c}{ENC} & \multicolumn{3}{c}{GC} & \multicolumn{3}{c}{GC3s} \\
    \cmidrule(lr){4-6} \cmidrule(lr){7-9} \cmidrule(lr){10-12}
    Genome &
        Species &
        \shortstack{Number \\ of Genes} &
        Mean & Median & StDev &
        Mean & Median & StDev &
        Mean & Median & StDev \\
    \midrule \startshading
    \texttt{BalGla} & \balgla{}                              & \num{16217} & 43.17 & 42.84 & 5.52 & 0.645 & 0.648 & 0.061 & 0.781 & 0.796 & 0.102 \\
    \texttt{BalCre} & \textit{Balanus crenatus}              & \num{7675}  & 42.23 & 41.40 & 4.78 & 0.645 & 0.649 & 0.055 & 0.815 & 0.837 & 0.093 \\
    \texttt{AmpAmp} & \textit{Amphibalanus amphitrite}       & \num{25242} & 44.33 & 43.93 & 5.05 & 0.632 & 0.636 & 0.054 & 0.775 & 0.792 & 0.099 \\
    \texttt{PolPol} & \textit{Pollicipes pollicipes}         & \num{17607} & 42.04 & 41.57 & 4.56 & 0.644 & 0.648 & 0.049 & 0.829 & 0.845 & 0.082 \\
    \texttt{DroMel} & \textit{Drosophila melanogaster}       & \num{12312} & 49.97 & 50.47 & 3.99 & 0.534 & 0.540 & 0.053 & 0.630 & 0.640 & 0.110 \\
    \texttt{AnoGam} & \textit{Anopheles gambiae}             & \num{11804} & 48.43 & 48.92 & 5.09 & 0.556 & 0.566 & 0.062 & 0.667 & 0.688 & 0.136 \\
    \texttt{CaeBre} & \textit{Caenorhabditis brenneri}       & \num{22904} & 50.95 & 51.55 & 3.21 & 0.436 & 0.430 & 0.045 & 0.411 & 0.399 & 0.088 \\
    \texttt{MagGig} & \textit{Magallana gigas}               & \num{24963} & 53.34 & 53.66 & 2.58 & 0.442 & 0.441 & 0.043 & 0.433 & 0.426 & 0.085 \\
    \texttt{StrPur} & \textit{Strongylocentrotus purpuratus} & \num{25759} & 53.92 & 54.39 & 2.68 & 0.481 & 0.480 & 0.035 & 0.489 & 0.486 & 0.079 \\
    \texttt{BraLan} & \textit{Branchiostoma lanceolatum}     & \num{23301} & 51.20 & 51.78 & 3.57 & 0.537 & 0.536 & 0.043 & 0.620 & 0.619 & 0.095 \\
    \texttt{CioInt} & \textit{Ciona intestinalis}            & \num{14370} & 52.49 & 52.75 & 2.66 & 0.424 & 0.422 & 0.037 & 0.369 & 0.367 & 0.069 \\
    \texttt{HomSap} & \textit{Homo sapiens}                  & \num{21573} & 49.39 & 50.57 & 4.51 & 0.538 & 0.541 & 0.087 & 0.593 & 0.603 & 0.163 \\
    \bottomrule
    \end{tabular}
\end{sidewaystable}

% Summary of ENC models
\begin{table}[htbp]
    \centering
    \small
    \suptablecaption{\textbf{Summary of ENC model fits across metazoan genomes.}
        \comment{update (remove?) this?}
        Table shows the fits for the model between the expected and observed effective
        number of codons (ENC) across 12 metazoan genomes.
        ENC deviation values represents the deviation from Wright's curve, as shown in
        Figure~\ref{fig:sup_enc_gc3s}.
        Larger ($\gg0$) values denote a reduction in the observed number of effective
        codons when compared to the GC3s-based expectation, suggesting translational
        selection.
        The linear regression values represent the linear relationship between the
        observed and expected number of effective codons, as seen in Figure~\ref{fig:sup_enc_reg}.
        Deviations in ENC product of translational selection are seen as a reduction in
        both slope and coefficient of determination ($\text{R}^2$).
        }
    \label{tab:sup_enc_models}
    \shadedrows
    \begin{tabular}{llrcccccccc}
    \toprule
    & & & \multicolumn{5}{c}{ENC deviation} & \multicolumn{3}{c}{Linear regression} \\
    \cmidrule(lr){4-8} \cmidrule(lr){9-11}
    Genome &
        Species &
        \shortstack{Number \\ of Genes} &
        Mean & Median & StDev & Min & Max &
        Slope & Intercept & $\text{R}^2$ \\
    \midrule \startshading
    \texttt{BalGla} & \bgla{}                  & \num{16217} & 0.065 & 0.053 & 0.090 & $-0.344$ & 0.571 & 0.580 & 16.18 & 0.492 \\
    \texttt{BalCre} & \textit{B. crenatus}     & \num{7675}  & 0.040 & 0.044 & 0.052 & $-0.235$ & 0.382 & 0.725 & 10.19 & 0.844 \\
    \texttt{AmpAmp} & \textit{A. amphitrite}   & \num{25242} & 0.052 & 0.055 & 0.047 & $-0.225$ & 0.517 & 0.747 & 9.27  & 0.868 \\
    \texttt{PolPol} & \textit{P. pollicipes}   & \num{17607} & 0.025 & 0.032 & 0.054 & $-0.265$ & 0.436 & 0.746 & 9.77  & 0.806 \\
    \texttt{DroMel} & \textit{D. melanogaster} & \num{12312} & 0.092 & 0.092 & 0.039 & $-0.125$ & 0.356 & 0.754 & 8.44  & 0.760 \\
    \texttt{AnoGam} & \textit{A. gambiae}      & \num{11804} & 0.074 & 0.075 & 0.040 & $-0.184$ & 0.480 & 0.779 & 7.58  & 0.871 \\
    \texttt{CaeBre} & \textit{C. brenneri}     & \num{22904} & 0.107 & 0.100 & 0.057 & $-0.079$ & 0.518 & 0.436 & 26.00 & 0.175 \\
    \texttt{MagGig} & \textit{M. gigas}        & \num{24963} & 0.079 & 0.076 & 0.040 & $-0.108$ & 0.567 & 0.560 & 20.89 & 0.330 \\
    \texttt{StrPur} & \textit{S. purpuratus}   & \num{25759} & 0.088 & 0.083 & 0.037 & $-0.287$ & 0.465 & 0.853 & 3.51  & 0.358 \\
    \texttt{BraLan} & \textit{B. lanceolatum}  & \num{23301} & 0.085 & 0.082 & 0.046 & $-0.166$ & 0.487 & 0.642 & 15.17 & 0.624 \\
    \texttt{CioInt} & \textit{C. intestinalis} & \num{14370} & 0.060 & 0.057 & 0.038 & $-0.156$ & 0.428 & 0.591 & 19.45 & 0.588 \\
    \texttt{HomSap} & \textit{H. sapiens}      & \num{21573} & 0.086 & 0.091 & 0.046 & $-0.201$ & 0.486 & 0.693 & 11.83 & 0.818 \\
    \bottomrule
    \end{tabular}
\end{table}


% Breakdown of the per-feature Pi
\begin{sidewaystable}[ht]
    \centering
    \small
    \setlength{\tabcolsep}{4pt}
    \suptablecaption{\textbf{Nucleotide diversity ($\pi$) across different genomic features.}
        For all features, pooled $\pi$ denotes the total sum of pairwise differences over
        the sum pairwise comparisons ($\sum\text{differences}/\sum\text{comparisons}$) for all
        accessible sites.
        Genome-wide and intergenic statistics are calculated over 10\,kbp windows with at
        \num{1000} accessible sites.
        Diversity over smaller genomic intervals (CDS, introns, 5' UTR, 3' UTR, and lncRNA)
        are calculated on windows with at least 50 accessible sites; $\geq20$ accessible sites
        for tRNAs.
        Diversity at 0-fold and 4-fold degenerate sites are calculated for exons with
        $\geq50$ accessible sites and pooled over the span of the whole gene.}
    \label{tab:sup_feature_pi}
    \shadedrows
    \begin{tabular}{lrrrrrrrrrrrr}
        \toprule
        & \multicolumn{4}{c}{$\pi$ ($\times10^{-2}$)} & &
            \multicolumn{3}{c}{Window length (bp)} & &
            \multicolumn{3}{c}{Accessible sites per window} \\
        \cmidrule(lr){2-5} \cmidrule(lr){7-9} \cmidrule(lr){11-13}
        Feature &
            Pooled & Mean & Median & StDev &
            \shortstack{Number \\ of \\ windows} &
            Mean & Median & StDev &
            \shortstack{Number of \\ accessible \\ sites} &
            Mean & Median & StDev \\
        \midrule \startshading
        % Feature                          pooledPi meanPi  medianPi sdPi    nWinds        meanWinLen       MedWinLen     sdWinLen         nBP               meanWinBP        medWinBP     sdWinBP
        \shortstack[l]{Genome\\wide}       & 2.45   & 2.32  & 2.35   & 1.27  & \num{63098} & \num{9999.18}  & \num{10000} & 65.59          & \num{222343515} & \num{3523.78}  & \num{3154} & \num{1862.62} \\
        Intergenic                         & 2.41   & 2.28  & 2.29   & 1.30  & \num{38889} & \num{9210.39}  & \num{10000} & \num{1961.53}  & \num{116883842} & \num{3005.58}  & \num{2514} & \num{1737.17} \\
        CDS                                & 1.61   & 1.52  & 1.35   & 1.12  & \num{77517} & 294.24         & 165         & 451.36         & \num{19880799}  & 256.47         & 158        & 376.96 \\
        Intron                             & 2.70   & 3.27  & 3.19   & 1.91  & \num{86036} & \num{3592.94}  & \num{1002}  & \num{8494.48}  & \num{94754847}  & \num{1101.34}  & 407        & \num{2741.38} \\
        5' UTR                             & 2.31   & 2.26  & 1.92   & 1.71  & \num{4636}  & 188.61         & 142         & 157.46         & \num{750453}    & 161.88         & 124        & 130.79 \\
        3' UTR                             & 2.36   & 2.41  & 2.19   & 1.60  & \num{8111}  & 670.82         & 559         & 485.01         & \num{4457431}   & 549.55         & 455        & 405.15 \\
        lncRNA                             & 1.98   & 1.76  & 1.52   & 1.32  & \num{26739} & 414.91         & 189         & 524.89         & \num{9385177}   & 350.99         & 176        & 441.44 \\
        tRNA                               & 0.470  & 0.481 & 0      & 1.02  & \num{435}   & 65.61          & 72          & 17.27          & \num{26650}     & 61.26          & 72         & 18.74 \\
        \shortstack[l]{0-fold\\degenerate} & 0.669  & 0.825 & 0.498  & 0.948 & \num{15647} & \num{1425.77}  & \num{1116}  & \num{1345.18}  & \num{12688329}  & 810.91         & 623        & 807.87 \\
        \shortstack[l]{4-fold\\degenerate} & 4.31   & 3.87  & 3.75   & 2.49  & \num{15647} & \num{1425.77}  & \num{1116}  & \num{1345.18}  & \num{3718088}   & 237.62         & 178        & 233.33 \\
    \bottomrule
    \end{tabular}
\end{sidewaystable}

% Genebank/RefSeq data used
\begin{table}[ht]
    \centering
    \suptablecaption{\textbf{Publicly available NCBI assemblies used for comparative
        analyses.} For the NCBI accessions, \texttt{GCF\_*} denotes RefSeq assemblies and
        annotations. \texttt{GCA\_*} accessions denote GenBank assemblies and (when available)
        annotations.}
    \shadedrows
    \begin{tabular}{lll}
        \toprule
            Genome & Species & NCBI accession \\
        \midrule \startshading
            \texttt{AmpAmp} & \textit{Amphibalanus amphitrite}       & \texttt{GCF\_019059575.1} \\
            \texttt{AmpAmp} & \textit{Amphibalanus amphitrite}       & \texttt{GCA\_037642225.1} \\
            \texttt{AmpImp} & \textit{Amphibalanus improvisus}       & \texttt{GCA\_964274985.1} \\
            \texttt{CapMit} & \textit{Capitulum mitella}             & \texttt{GCA\_030062745.1} \\
            \texttt{PolPol} & \textit{Pollicipes pollicipes}         & \texttt{GCF\_011947565.3} \\
            \texttt{DroMel} & \textit{Drosophila melanogaster}       & \texttt{GCF\_000001215.4} \\
            \texttt{AnoGam} & \textit{Anopheles gambiae}             & \texttt{GCF\_943734735.2} \\
            \texttt{CaeBre} & \textit{Caenorhabditis brenneri}       & \texttt{GCA\_964036135.1} \\
            \texttt{MagGig} & \textit{Magallana gigas}               & \texttt{GCF\_963853765.1} \\
            \texttt{StrPur} & \textit{Strongylocentrotus purpuratus} & \texttt{GCF\_000002235.5} \\
            \texttt{BraLan} & \textit{Branchiostoma lanceolatum}     & \texttt{GCF\_035083965.1} \\
            \texttt{CioInt} & \textit{Ciona intestinalis}            & \texttt{GCF\_018327825.1} \\
            \texttt{HomSap} & \textit{Homo sapiens}                  & \texttt{GCF\_000001405.40} \\
        \bottomrule
    \end{tabular}
    \label{tab:sup_ncbi_genomes}
\end{table}

% SRA data
\begin{table}[ht]
    \centering
    \suptablecaption{\textbf{Publicly available NCBI short-read archive data used
        for comparative analyses.}}
    \shadedrows
    \begin{tabular}{lll}
        \toprule
            SRA & BioProject & Species \\
        \midrule \startshading
            % SRA                  BioProject             Species
            \texttt{SRR8775109}  & \texttt{PRJNA528777} & \textit{Amphibalanus improvisus} \\
            \texttt{SRR8775110}  & \texttt{PRJNA528777} & \textit{Amphibalanus improvisus} \\
            \texttt{SRR8775111}  & \texttt{PRJNA528777} & \textit{Amphibalanus improvisus} \\
            \texttt{SRR8775112}  & \texttt{PRJNA528777} & \textit{Amphibalanus improvisus} \\
            \texttt{SRR8775113}  & \texttt{PRJNA528777} & \textit{Amphibalanus improvisus} \\
            \texttt{SRR8775114}  & \texttt{PRJNA528777} & \textit{Amphibalanus improvisus} \\
            \texttt{SRR8775115}  & \texttt{PRJNA528777} & \textit{Amphibalanus improvisus} \\
            \texttt{SRR8775116}  & \texttt{PRJNA528777} & \textit{Amphibalanus improvisus} \\
            \texttt{SRR8775117}  & \texttt{PRJNA528777} & \textit{Amphibalanus improvisus} \\
            \texttt{SRR8775118}  & \texttt{PRJNA528777} & \textit{Amphibalanus improvisus} \\
            \texttt{SRR8775119}  & \texttt{PRJNA528777} & \textit{Amphibalanus improvisus} \\
            \texttt{SRR8775120}  & \texttt{PRJNA528777} & \textit{Amphibalanus improvisus} \\
            \texttt{SRR8775121}  & \texttt{PRJNA528777} & \textit{Amphibalanus improvisus} \\
            \texttt{SRR8775122}  & \texttt{PRJNA528777} & \textit{Amphibalanus improvisus} \\
            \texttt{SRR8775123}  & \texttt{PRJNA528777} & \textit{Amphibalanus improvisus} \\
            \texttt{SRR8775124}  & \texttt{PRJNA528777} & \textit{Amphibalanus improvisus} \\
            \texttt{SRR8775125}  & \texttt{PRJNA528777} & \textit{Amphibalanus improvisus} \\
            \texttt{SRR8775126}  & \texttt{PRJNA528777} & \textit{Amphibalanus improvisus} \\
            \texttt{SRR8775127}  & \texttt{PRJNA528777} & \textit{Amphibalanus improvisus} \\
            \texttt{SRR8775128}  & \texttt{PRJNA528777} & \textit{Amphibalanus improvisus} \\
            \texttt{SRR8775129}  & \texttt{PRJNA528777} & \textit{Amphibalanus improvisus} \\
            \texttt{SRR8775130}  & \texttt{PRJNA528777} & \textit{Amphibalanus improvisus} \\
            \texttt{SRR8775131}  & \texttt{PRJNA528777} & \textit{Amphibalanus improvisus} \\
            \texttt{SRR8755479}  & \texttt{PRJNA528169} & \textit{Amphibalanus improvisus} \\
            \texttt{SRR18959849} & \texttt{PRJNA816681} & \textit{Capitulum mitella} \\
            \texttt{SRR14354747} & \texttt{PRJNA725059} & \textit{Capitulum mitella} \\
            \texttt{SRR20067441} & \texttt{PRJNA856037} & \textit{Capitulum mitella} \\
            \texttt{SRR20067442} & \texttt{PRJNA856037} & \textit{Capitulum mitella} \\
            \texttt{SRR20067443} & \texttt{PRJNA856037} & \textit{Capitulum mitella} \\
            \texttt{SRR20067446} & \texttt{PRJNA856037} & \textit{Capitulum mitella} \\
            \texttt{SRR20067447} & \texttt{PRJNA856037} & \textit{Capitulum mitella} \\
            \texttt{SRR20067448} & \texttt{PRJNA856037} & \textit{Capitulum mitella} \\
            \texttt{SRR20067449} & \texttt{PRJNA856037} & \textit{Capitulum mitella} \\
            \texttt{SRR20067444} & \texttt{PRJNA856037} & \textit{Capitulum mitella} \\
            \texttt{SRR20067445} & \texttt{PRJNA856037} & \textit{Capitulum mitella} \\
            \texttt{SRR13527601} & \texttt{PRJNA681321} & \textit{Capitulum mitella} \\
            \texttt{SRR13527602} & \texttt{PRJNA681321} & \textit{Capitulum mitella} \\
            \texttt{SRR13527603} & \texttt{PRJNA681321} & \textit{Capitulum mitella} \\
            \texttt{SRR13527604} & \texttt{PRJNA681321} & \textit{Capitulum mitella} \\
        \bottomrule
    \end{tabular}
    \label{tab:sup_sra_data}
\end{table}

